% ======================================================================
% v7 DRAFT: abstract + Section 1 + Section 2, rewritten from the reframe
% (motivation that keeps a paradigm is a commitment to a generative model,
% not a preference over its conclusions). Drop-in replacement for the
% corresponding blocks of changing_networked_mind_v6.tex; everything from
% \section{The model} onward is unchanged.
%
% Each paragraph is tagged: [KEPT] = v6 text, at most lightly edited;
% [REFRAMED] = v6 material, reordered and rewritten; [NEW] = not in v6.
% Two new bib keys at the bottom (nielsen2022, quine1978): VERIFY before use.
% Bib keys feyerabend1975, lehrer1981, laudan1977 already exist in v6 and
% were never cited in its text; they are cited here.
% ======================================================================

\begin{abstract}
% [REFRAMED] Opens with the question rather than the Kuhn thesis; keeps the
% core/belt-from-geometry sentence, the contraction sentence and the closing
% sentence; adds the tilt finding and the proliferation line.
To change a paradigm is to change the frame in which beliefs are posed, not a credence
within it. Peripheral commitments drift under accumulating evidence while central ones
resist, and when they give way they give way quickly. We ask what lets a community of
bounded reasoners hold a structure of commitments past disconfirming evidence, and what
finally makes the holding give way. We model this within active inference, building on a
recent account of belief updating as motivated inference subject to a cost of changing
one's mind. Each agent's commitments form a directed web of belief in which the cost of
revising a commitment is the surrounding inference that must re-equilibrate when it moves,
so Lakatos's hard core and protective belt fall out of the precision geometry rather than
being decreed. Agents exchange beliefs, and pieces of structure, over a communication
network. Characterising the dynamics analytically and in simulation, we find two things.
First, the field's default aggregation step, precision-weighted averaging, cannot reproduce
the signature of paradigm change: it contracts the community onto a compromise no member
holds, whatever the parameters. Second, the familiar explanation of resistance, attachment
to the conclusions one wants true, contributes only timing: entered as a value that tilts
inference, it shifts when a transition arrives, not whether. What holds a paradigm is a
commitment to a structure. Giving each agent a posterior over rival wirings, each anchored
to its own prior, restores the capacity the default lacks and recovers the Kuhn cycle with
crisis as a measurable state of belief: normal science, crisis, revolution, new normal.
This is Feyerabend's principle of proliferation made computational: without a represented
rival, an anomaly cannot register as crisis. We close by discussing implications for
engineered communities of reasoning agents, whose aggregation protocols decide the same
question by design.

\keywords{Active inference \and Structure learning \and Belief sharing \and Paradigm change \and Social epistemology}
\end{abstract}

% ======================================================================
\section{Introduction: changing a frame, not a credence}
\label{sec:background}
% ======================================================================

% [REFRAMED] v6's second paragraph, promoted to open the paper; the three
% philosophers from v6's first sentence are folded in with one clause each.
% "What his account leaves open is a mechanism" is gone (reviewer 53cW).
To change one's mind about a fact is to update a belief; to change a paradigm is to change
the \emph{frame} in which beliefs are posed: not a credence but a \emph{structure}, a
network of conditional dependencies through which information propagates. When the chemical
community of the late eighteenth century gave up phlogiston, it did not revise a single
credence: it rewrote a network in which combustion, calcination, respiration, and the theory
of acids were all load-bearing on a hub that had to go
\citep{blumenthal2017,holmes2000,boantza2011}. Phlogiston was load-bearing for too much.
Kuhn described this cycle historically \citep{kuhn1962}; Lakatos gave the held structure a
shape, a hard core and a protective belt \citep{lakatos1970}; Quine gave it a name, a web of
belief that meets experience only as a whole \citep{quine1951}. This paper asks what it
takes for a community of bounded reasoners to run that cycle: which representation lets it
hold a structure of commitments past the first disconfirming evidence, and what finally
makes the holding give way.

% [NEW] The attachment paragraph: the tempting answer, built in, found wanting.
The familiar answer is attachment. A result on which a career rests is held more tightly
than its evidence warrants \citep{kitcher1990,strevens2003}; a paradigm, on this view, is
held by wanting it true. We build that answer into the model, as an intrinsic value on each
commitment that propagates through the web and tilts inference toward preferred conclusions
\citep{hyland2025}, and we find that its measured contribution is timing: it shifts when a
transition arrives, not whether. What holds a paradigm in the model is something else. It is
a generative model of one's own to reason from, and a rule for sharing beliefs that does not
average that model away. Attachment to conclusions is not enough; commitment to a structure
is what survives contact with a rival. Philosophy of science has long drawn the distinction
in words, between preferring a theory and seeing the world through it
\citep{kuhn1962,feyerabend1975}; the model draws it in mechanism, and finds only the second
causally effective.

% [REFRAMED] v6's third paragraph (network epistemology, ECHO, active inference),
% now opened by the Feyerabend ancestor and the three named couplings; ECHO's
% shortfall is spelled out (reviewer HadU); Lehrer & Wagner cited as the default.
The result has an ancestor. Feyerabend argued that the evidence against a theory often
cannot be seen until an alternative is on the table: it takes a rival to turn a puzzle into
an anomaly \citep{feyerabend1975}. In the model this is literal. Without a represented rival
the community has no belief state that means \emph{torn}, and the world's turn arrives as
growing confidence in a compromise no one holds. The rest of the model's vocabulary is drawn
from the same tradition, and we say so as we build it: Quine's maxim of minimum mutilation
is the cost agents pay to revise their web \citep{quine1978}; Lakatos's negative heuristic,
reason from the core and let the evidence land in the belt, is the anchoring that keeps a
candidate alive \citep{lakatos1970}; and the sharing rule we begin from is the rational
consensus of Lehrer and Wagner \citep{lehrer1981,degroot1974}, so the contraction result
below is a result about that tradition. The social side belongs to formal social
epistemology, which asks when a community converges too soon
\citep{zollman2010,zollman2013,oconnor2019,jonard2025,michelini2023}. That literature, in
its standard form, updates a scalar credence in a fixed proposition on a trust graph. It can
say that a community herds on a number or splits over it; it cannot say how revising one
commitment forces a re-reading of the rest, because its agents cannot hold different kinds
of world. Thagard's ECHO \citep{thagard1989,thagard1990} comes closest, representing a
theory as a coherence network, but its dependencies are undirected, it carries no cost of
changing one's mind, and it models one reasoner rather than a community. Active inference
supplies the missing pieces one at a time: precision-based lock-in \citep{albarracin2022},
warnings that naive aggregation degrades collective precision
\citep{friston2023federated,waade2025}, and a single-agent cost of mind-change
\citep{hyland2025}. What none of them does is make belief structure itself the object of
inference within a community.

% [KEPT] v6's second-audience material, from its first paragraph, with the
% metascience citation added.
The question has a second audience. Communities of reasoning agents are now built as well
as studied, and every such community commits, in its aggregation step, to an answer it
rarely examines: what happens to the member who disagrees. The social processes of science
are themselves increasingly treated as designable \citep{nielsen2022}. Machine learning
names the property at stake \emph{open-endedness} \citep{stanley2017oe,hughes2024}; social
epistemology catalogues its failure modes \citep{nguyen2018,zollman2010}. What neither
supplies is a foundation: a model of a community of structure-learning agents explicit
enough that one can prove which aggregation dynamics preserve open-endedness and which
foreclose it. The problem, throughout, is averaging: taking the mean refines the frame
everyone shares, and can never grow a new one.

% [REFRAMED] v6's closing paragraph; adds the "names the idea it instantiates"
% promise for Section 3 and the three-homes-of-motivation sentence.
This paper lets a piece of structure travel, so that one agent can come to represent a
commitment its peers do not. Its agents carry a directed network of commitments and revise
it by Bayesian model reduction; expansion, the conceiving of a commitment no one yet holds,
is deferred by scope (\S\ref{sec:structurelearning}), but when agents are given new
structures they can take them up. They are bounded, sampling a complex world through narrow
channels, valuing some conclusions over others, and judging what they receive against their
own model. Section~\ref{sec:bmebmr} assembles this community from the field's own default
choices and names, for each choice, the idea from philosophy of science it instantiates.
Section~\ref{sec:results} shows what those defaults force: pooling is a contraction, so a
conflicting source is never rejected, only averaged in; and what undoes it: agents that
represent rival structures, each anchored to its own prior, keep different readings of the
same world apart and run the Kuhn cycle with crisis as a state of the belief. The same
experiments let us read off the three places motivation can enter a bounded reasoner, the
mean, the precision, and the model index, and only the last holds. The discussion draws the
map: what filling the representational slot bought, and what remains.

% ======================================================================
\section{Desiderata for a structural model of science}
\label{sec:hidden}
% ======================================================================

% [KEPT] First ingredient unchanged.
A structural model of science needs at least three ingredients, and each has a default
formalization so natural it is hardly ever examined. First, structure as the object of
inference: a scientist confronting a new domain is not given the right variables, let alone
the dependencies among them; frames differ in kind, not in degree, so credences over a fixed
menu cannot carry them \citep{ullman2020,kemp2008form,ullman2012}. Active inference can
model this: when structure is the object of inference, exploration need not be added by
hand, as the $\epsilon$ of an $\epsilon$-greedy algorithm is \citep{zollman2010,oconnor2019};
it can be cast as the epistemic value of expanding the model to entertain a commitment not
yet in it \citep{friston2015,friston2017curiosity}, with Bayesian model reduction as its
complement, pruning what the data abandon in closed form \citep{friston2018bmr,neacsu2022}.
To model structure learning, then, is to model how a Bayes net learns to change itself over
time. The default is a precision-weighted Gaussian model over the network.

% [REFRAMED] The value ingredient now states the three levels; the default is
% named as the first of them.
Second, \emph{value}, or motivated belief updating: the commitments a scientist already has
about what they care about. A result on which a career rests is held more tightly than its
evidence warrants \citep{kitcher1990,strevens2003}, and attitudes demonstrably live on
networks of exactly this kind \citep{dalege2016}. Value can enter a bounded reasoner at three
levels: as a preference over conclusions, tilting the mean of a belief; as attention,
deciding which evidence is weighed; or as a commitment to a structure, deciding which model
one reasons from. The default puts it at the first level: a utility that tilts the prior
\citep{hyland2025}. Kuhn's values for theory choice, simplicity, scope, consistency,
fruitfulness, live at the third \citep{kuhn1977}; part of what the experiments below
measure is the distance between the two.

% [KEPT] Third ingredient, with Lehrer & Wagner added to the default.
Third, communication, because science is social: beliefs travel a trust graph and must
somehow be combined when they meet. The default is precision-weighted averaging, the
information form of classical opinion pooling and of rational consensus
\citep{degroot1974,genest1986,lehrer1981}. Section~\ref{sec:bmebmr} builds the community
out of exactly these defaults, on purpose. The question the rest of the paper asks is
whether such a community can ever do more than refine: whether agents built this way can
grow a frame, and not merely polish the one they hold; and, where they cannot, which of the
three defaults has to change.

% ======================================================================
% NEW BIB ENTRIES (append inside \begin{thebibliography}); VERIFY all three.
% ======================================================================
% \bibitem[Nielsen and Qiu(2022)]{nielsen2022}
% M.~Nielsen and K.~Qiu.
% A vision of metascience: an engine of improvement for the social processes of science.
% Essay, 2022. https://scienceplusplus.org/metascience/
%
% \bibitem[Quine and Ullian(1978)]{quine1978}
% W.~V.~O.~Quine and J.~S.~Ullian.
% \emph{The Web of Belief}, 2nd edition. Random House, 1978.
% (The maxim of minimum mutilation; also stated in Quine, \emph{Philosophy of Logic}, 1970.)
%
% \bibitem[Kuhn(1977)]{kuhn1977}
% T.~S.~Kuhn.
% Objectivity, value judgment, and theory choice.
% In \emph{The Essential Tension}, pages 320--339. University of Chicago Press, 1977.
